% ======================================================================
% 前置部分：前言、符号约定、编者前言
% （根据 1995 年 Westview Press 版转写并翻译）
% ======================================================================
\chapter*{前言}
\addcontentsline{toc}{chapter}{前言}

量子场论是一套将现代物理学的三大主题——量子理论、场概念与相对性原理——结合起来的观念和工具。如今，大多数从事实际工作的物理学家都需要了解一定的量子场论，还有许多人对此感到好奇。该理论是现代基本粒子物理的基础，并为核物理、原子物理、凝聚态物理和天体物理提供了必不可少的工具。此外，量子场论还在物理学与数学之间架起了新的桥梁。

人们也许会认为，这样一个威力强大且应用广泛的理论必定复杂而艰深。事实上，量子场论的核心概念与技术相当简单且直观。对于那些量子场论研究者经常使用的各种形象化工具（费曼图、重整化群流、对称变换的空间）而言尤其如此。诚然，掌握这些工具需要时间，而用严格的证明把整个理论体系联系起来则可能变得极其技术化。尽管如此，我们仍认为，量子场论的基本概念和工具可以做到让所有物理学家都能掌握，而不仅仅是少数精英专家。

此前已有若干著作成功地让这一理论的某些部分变得易于学生理解。其中最负盛名的，当属 20 世纪 60 年代我们的斯坦福同事 Bjorken 和 Drell 所著的两卷本教科书。在我们看来，他们的著作把抽象的形式体系、直观的解释和实际的计算完美地融合在一起，且都以极大的细致与清晰呈现。然而，自 20 世纪 60 年代以来，量子场论无论在概念框架（重整化群、新型对称性）还是应用领域（凝聚态体系中的临界指数、基本粒子物理的标准模型）上都得到了巨大的发展。一本既能对这门学科（包括上述新进展）作出全面综述、又能提供与 Bjorken 和 Drell 同样易于理解且论述同样深入的量子场论教科书，早该问世了。我们正是怀着这一目标来写作本书的。

\section*{全书概览}

本教科书由三大部分组成。第一部分主要讨论电磁学的量子理论，它为具有直接实验应用的量子场论提供了第一个范例。第三部分主要讨论出现在粒子相互作用标准模型中的那些具体的量子场论。第二部分则是这两个主题之间的桥梁；它意在尽可能直截了当地引入量子场论中一些非常深刻的概念。

第一~部分从研究满足线性运动方程的场——即没有相互作用的场——开始。在这里我们探讨量子力学与狭义相对论相结合的含义，并了解粒子是如何作为场的量子化激发而产生的。随后我们引入这些粒子之间的相互作用，并发展出一套系统的方法来计算它们的影响。在完成这一引入之后，我们在电磁学量子理论中进行具体的计算。这些计算既展示了电子和光子行为的特殊之处，也展示了相互作用量子场行为的一些普遍特征。

在第一~部分的若干计算中，朴素的处理方法会导致无穷大的结果。无穷大的出现是量子场论中一个众所周知的特征。它有时被当作量子场论不自洽的证据（尽管对点粒子的经典电动力学也可以提出类似的论证）。很长一段时间里，人们认为只要这样来组织计算——使得可以直接与实验比较的量中不出现无穷大——就足够了。然而，近期发展的一个重要洞见是，这些形式上的无穷大实际上包含着可用于预言体系定性行为的重要信息。在本书第二~部分，我们系统地发展这一关于无穷大的理论。这一发展利用了量子力学涨落与热涨落之间的类比，从而使量子场论与统计力学之间建立起一座桥梁。在第二~部分的末尾，我们讨论量子场论在凝聚态体系相变理论中的应用。

第三~部分讨论量子电动力学的各种推广，它们带来了描述基本粒子之间作用力的成功模型。为了推导这些推广，我们首先分析并推广电动力学的基本对称性，然后讨论对这一具有推广对称性的理论进行量子化所带来的后果。这一分析导致前面引入的概念得到复杂而相当非平凡的应用。我们以粒子物理标准模型的介绍及其若干实验检验的讨论来结束第三~部分。

本书的结语定性地讨论了量子场论研究的前沿领域，并给出了能够引导学生进入下一学习阶段的参考文献。

在观点可以选择之处，我们通常选用适合基本粒子物理应用的表述方式来解释各种思想。这一选择反映了我们的背景和研究兴趣。它也反映了我们一个坚定的信念——即便在当今这个思想相对主义的时代，在最深层次上揭示大自然的行为仍是一件非同寻常的事。我们为能以基本相互作用的结构作为研究对象而感到自豪，并希望能把这一追求的宏伟气势和持久的生命力传达给读者。

\section*{如何使用本书}

本书是对量子场论的一篇导论。这里我们首先且最重要的是要说明：我们假定读者先前并不了解这门学科。本书的水平应适合首次学习量子场论课程的学生，在美国这类课程通常在研究生二年级开设。我们假定学生已经完成了经典力学、经典电动力学和量子力学的研究生课程。在第二~部分我们还假定学生具备一定的统计力学知识。不过，并不需要掌握这些课程所涵盖的每一个专题。至关重要的先修内容包括：动力学的拉格朗日表述和哈密顿表述、用张量记号写出的电磁学相对论表述、用阶梯算符对谐振子的量子化，以及非相对论量子力学中的散射理论。数学上的先修要求包括：对旋转群在自旋量子力学中应用的理解，以及一定的复平面围道积分能力。

尽管自称是一篇``导论''，本书仍相当冗长。在某种程度上，这是因为书中包含大量的具体计算和已演算的例题。然而，我们必须承认，所涵盖的专题总数也相当多。即使是专攻基本粒子理论的学生也会发现，他们的第一个研究课题只需用到本书的一部分内容，还需要从研究文献中搜集额外的专门专题。尽管如此，我们仍然认为，那些想要成为基本粒子理论专家、并充分理解其对基本相互作用的统一观点的学生，应当掌握本书的每一个专题。而主要兴趣在其他物理领域或粒子实验的学生，则可以选择一条简短得多的``导论''路线，略去若干章。

本书的第一作者曾在斯坦福大学为期一年的课程中成功``讲完''了本书 90\% 的内容。但这是一个错误；以这样的进度，对于基础一般的学生来说，没有足够的时间来消化这些内容。我们较为明智的同事们发现，一学期讲授本书大约一个部分更为合理。因此，在规划为期一年的量子场论课程时，他们要么把第三~部分留待更高级的学习阶段，要么从第二~部分和第三~部分中选取约一半的内容，其余部分留给学生自学。

我们把本书设计成可以逐页通读，系统地引入该领域的所有主要思想。或者，读者也可以走一条快捷路线，它强调对基本粒子物理中较不形式化的应用，足以让学生为该领域的实验或唯象研究做好准备。可以从这条快捷路线中略去的章节在目录中都用星号标出；未标星号的章节都不依赖于这些较深入的内容。在未标星号的章节中，次序也可以作一定调整：第~10 章不依赖于第~8 章和第~9 章；第~11.1 节直到第~20 章之前才需要；第~20 章和第~21 章与第~17 章相互独立。

希望研读部分（而非全部）较深入章节的读者，应注意下面这张依赖关系表：
\begin{center}
\begin{tabular}{ll}
\toprule
在阅读 \ldots 之前 & 应先通读 \ldots\\
\midrule
第 13 章            & 第 11、12 章\\
第 16.6 节          & 第 11 章\\
第 18 章            & 第 12.4、12.5、16.5 节\\
第 19 章            & 第 9.6、15.3 节\\
第 19.5 节          & 第 16.6 节\\
第 20.3 节          & 第 19.1--19.4 节\\
第 21.3 节          & 第 11 章\\
\bottomrule
\end{tabular}
\end{center}

在每一章内部，标有星号的章节应按顺序阅读，唯一的例外是第~16.5 节和第~16.6 节并不依赖于第~16.4 节。

主要兴趣在统计力学的学生，可能会希望按顺序通读本书，直面第二~部分的深层形式问题，而略过大部分对高能现象才有意义的内容——即第三~部分的大部分。（不过，第~15 章、第~19 章以及第~20.1 节的内容，在凝聚态物理中确实有精彩的应用。）

我们向所有学生强调，在学习过程中主动地运用这些材料十分重要。不仔细推演每一个推导的中间步骤，恐怕不可能真正理解本书的任何一节。此外，每章末尾的习题阐明了这些一般思想，并常常把它们应用于非平凡的现实情境之中。然而，量子场论中最具启发性的练习对于普通的课外作业来说太长了，它们更接近小型研究课题的规模。我们已经在本书三个部分的末尾各提供了一道这样的长习题，并把它分解成带有提示和指导的小段。这些习题所需的时间与纸墨，将会是值得的投入。

在每个部分的开始，我们都加入了一章简短的``导引''，预先介绍一些即将讨论的思想和应用。由于这些章节比本书其余部分稍为容易，我们力劝所有学生都去阅读它们。

\section*{本书不是什么}

尽管我们希望本书能够为量子场论提供扎实的基础训练，但它绝不是一种完整的教育。有志于物理的学生将会希望在许多方面补充我们的论述。我们在此概括其中最重要的几点。

首先，这是一本关于理论方法的书，而不是对观测现象的综述。我们既不综述那些导致基本粒子物理标准模型的关键实验，也不详细讨论后来证实其预言的实验。类似地，在处理统计力学应用的章节中，我们不讨论那些导致了场论模型获得确认的丰富多彩的相变实验。我们强烈建议学生在阅读本书的同时，并行阅读对这些领域实验发展的现代阐述。

虽然我们详尽地阐述了量子场论的基本方面，但也会不加证明地陈述一些较深入的结果。例如，严格说来，在量子电动力学标准展开的所有阶，形式上的无穷大都可以从一切实验预言中消除。这一结果即所谓的\emph{可重整化性}，它具有重要的后果，我们将在第二~部分进行探讨。我们不给出可重整化性的一般证明。然而，我们确实在具有说明性的低阶计算中明确地演示了可重整化性，直觉地讨论了完整证明中出现的各种问题，并给出更完整证明的参考文献。更一般地说，我们力图把最重要的结果讲清楚（通常借助具体的例子），同时略去冗长而纯属技术性的推导。

任何入门性的综述都必须把某些专题划归到其范围之外。我们的方针是：纳入那些通过考虑弱相互作用的粒子和场、并利用相互作用强度作级数展开所能学到的量子场论内容。用这种方式竟能获得如此之多的洞见，实在令人惊叹。然而，这样界定我们的学科，便略去了束缚态理论，也略去了与非线性场方程非平凡解相关的各种现象。在结语中，我们给出了这类深入专题的更完整的清单。

最后，本书并未试图对量子场论的历史给出精确的记录。物理专业的学生确实需要了解物理学史，原因有很多。最重要的原因，是要获得对这门学科实验基础的精确理解。第二个重要原因，是要了解科学作为一项人类事业是如何进步的，了解个人的一个个小步骤是如何汇成整个科学共同体重大成就的。%
\footnote{量子场论与粒子物理的发展历史，近来已在一系列会议论文集中得到综述与讨论：\emph{The Birth of Particle Physics}（《粒子物理的诞生》），L.~M.~Brown 与 L.~Hoddeson 编（剑桥大学出版社，1983）；\emph{Pions to Quarks}（《从$\pi$介子到夸克》），L.~M.~Brown、M.~Dresden 与 L.~Hoddeson 编（剑桥大学出版社，1989）；以及 \emph{The Rise of the Standard Model}（《标准模型的兴起》），L.~M.~Brown、M.~Dresden、L.~Hoddeson 与 M.~Riordan 编（剑桥大学出版社，1995）。量子电动力学的早期历史在 Schweber（1994）那本引人入胜的书中得到了叙述。}

在本书中，这两种需求我们都没有去满足。相反，我们只包含了那种神话式的历史——其目的在于引出新思想并给它们命名。一条物理学原理通常有一个名字，而这个名字是按科学界就谁应当因其发展而获得荣誉所形成的共识来命名的。通常实际的功绩只占一部分，而真实的历史发展过程则相当复杂。但如果物理学家之间要相互交流，明确地命名就是必不可少的。

这里有一个例子。在第~17.5 节中，我们讨论一组支配质子结构的方程，它们通常被称为 Altarelli--Parisi 方程。我们的推导采用了 Gribov 和 Lipatov（GL）的方法。GL 的原始结果后来由 Christ、Hasslacher 和 Mueller（CHM）用一种更抽象的语言重新导出。在发现了强相互作用的正确基本理论（QCD）之后，Georgi、Politzer、Gross 和 Wilczek（GPGW）利用 CHM 的技术导出了描述质子结构变化的若干形式方程。Parisi 给出了若干独立推导中的第一个，把这些方程化成了有用的形式。他的工作与 GPGW 工作的结合，便给出了我们在第~18.5 节中所呈现的那些方程的推导。Dokshitzer 后来直接运用 GL 的方法，更简单地得到了这些方程。一段时间之后，Altarelli 和 Parisi 又独立地、循着同样的途径再次得到了这些方程。这两位作者还普及了这一技术：他们把它解释得非常清楚，鼓励实验物理学家用这些方程来解释他们的数据，并推动理论家计算这一图景的系统高阶修正。在第~17.5 节中，我们呈现了通往这条曲折历史之路终点的最短路径，并把``Altarelli--Parisi''这个名称加在了最终结果上。

学生阅读物理学史还有第四个理由：最初的突破性论文虽然缺乏教科书那种事后之明的优势，却常常充满精彩卓绝的个人洞见。我们强烈鼓励学生尽可能回到原始文献，去看看这一领域的创立者们当时在想什么。我们力图通过脚注中提供的参考文献来帮助这样的学生。虽然偶尔我们会仅为肯定功绩而引用某些论文，但大多数参考文献之所以被列入，是因为我们认为读者不应错过这些论文作者所提出的独特见解。

\section*{致谢}

如果没有我们众多师长、同事和朋友的帮助与鼓励，本书的写作是不可能完成的。我们非常高兴向这些人表示感谢。

Michael 有幸向这门学科当代的三位大师——Sidney Coleman、Steven Weinberg 和 Kenneth Wilson——学习场论。他还要感谢许多其他师长，包括 Sidney Drell、Michael Fisher、Kurt Gottfried、John Kogut 和 Howard Georgi，以及众多合作者，特别是 Orlando Alvarez、John Preskill 和 Edward Witten。他与康奈尔和斯坦福高能物理实验室的联系，以及与 Gary Feldman、Martin Perl、Morris Swartz 等人的讨论，也塑造了他的观点。Michael 向这些人，以及多年来在许多物理问题上给予他指导的许多人，表达他的感激之情。

Dan 从 Savas Dimopoulos、Leonard Susskind、Richard Blankenbecler 以及许多其他师长和同事那里学到了场论，他向他们表示感谢。对于他更广博的物理学教育，他要感谢许多老师，尤其是 Thomas Moore。此外，他感谢所有多年来对他的写作提出批评的老师与朋友。

对于在本书的构思与写作过程中给予帮助的人，许多人都应该得到更具体的感谢。本书起源于在康奈尔和斯坦福讲授的研究生课程，我们感谢这些课程的学生们的耐心、批评与建议。在本书写作的那些年里，我们得到了斯坦福直线加速器中心和（美国）能源部的支持，Dan 还得到了皮尤慈善信托基金会以及波莫纳学院、格林内尔学院和韦伯州立大学物理系的支持。我们特别感谢 Sidney Drell 的支持与鼓励。Michael 在加利福尼亚大学圣克鲁兹分校度过的学术休假极大地促进了本书的完成，他感谢那里的 Michael Nauenberg 以及物理系师生们的热情款待。

我们广泛传阅了本书的初稿，并一直受到对初稿的回应以及由此引发的大量建议和更正所鼓舞。以这种方式帮助我们的人包括 Curtis Callan、Edward Farhi、Jonathan Feng、Donald Finnell、Paul Frampton、Jeffrey Goldstone、Brian Hatfield、Barry Holstein、Gregory Kilcup、Donald Knuth、Barak Kol、Phillip Nelson、Matthias Neubert、Yossi Nir、Arvind Rajaraman、Pierre Ramond、Katsumi Tanaka、Xerxes Tata、Arkady Vainshtein、Rudolf von B\"unau、Shimon Yankielowicz 和 Eran Yehudai。对于更具体的技术问题或疑问方面的帮助，我们感谢 Michael Barnett、Marty Breidenbach、Don Groom、Toichiro Kinoshita、Paul Langacker、Wu-Ki Tung、Peter Zerwas 和 Jean Zinn-Justin。我们感谢我们的斯坦福同事 James Bjorken、Stanley Brodsky、Lance Dixon、Bryan Lynn、Helen Quinn、Leonard Susskind 和 Marvin Weinstein，书中的许多讲解都在他们身上做过检验。我们特别感谢 Michael Dine、Martin Einhorn、Howard Haber 和 Renata Kallosh，以及他们那些饱受折磨的学生——在数年时间里，他们一直把本书的初稿当作教材使用，并给我们带回了重要而有启发性的批评。

我们在此停下来，缅怀本书的两位未能看到它完成的朋友——Donald Yennie，他漫长的学术生涯使他成为量子电动力学的英雄人物之一；以及 Brian Warr，他在理论物理方面辉煌的潜质被艾滋病过早地扼杀了。

我们感谢我们在艾迪生-韦斯利的编辑——Allan Wylde、Barbara Holland、Jack Repcheck、Heather Mimnaugh 和 Jeffrey Robbins——感谢他们对这一项目的鼓励，以及始终没有失去对本书终将完成的信心。我们感谢 Mona Zeftel 和 Lynne Reed 在本书排版格式方面的建议，感谢 Jeffrey Towson 和 Cristine Jennings 为插图所做的娴熟准备工作。最终的排版页面是在 WSU（韦伯州立大学）印刷服务部制作的，并得到了 Allan Davis 的大力帮助。当然，如果没有许多让我们的计算机和网络连接保持正常运转的人，本书的制作是不可能完成的。

我们感谢在这漫长的工作过程中给予我们支持、使我们得以依赖的许多朋友。我们尤其感谢我们的父母：Dorothy 和 Vernon Schroeder，以及 Gerald 和 Pearl Peskin。Michael 还要感谢 Valerie、Thomas 和 Laura 的爱与理解。

最后，我们感谢你们——读者——为研读本书所付出的时间和精力。尽管我们力图清除本书中的概念性错误和印刷错误，但我们仍要预先为那些漏网之鱼表示歉意。我们将乐于听取你们对进一步改进量子场论表述方式的意见与建议。

\vspace{2em}
\noindent Michael E. Peskin\\
Daniel V. Schroeder

% ======================================================================
\chapter*{符号约定}
\addcontentsline{toc}{chapter}{符号约定}

\section*{单位制}

我们将使用``天赐''单位制，其中
\begin{equation}
  \hbar = c = 1 .
\end{equation}
因此，粒子的质量 $m$ 等于其静能 $mc^2$，也等于其康普顿波长的倒数 $mc/\hbar$。例如，
\begin{equation}
  m_\mathrm{electron} = 9.109 \times 10^{-28}\,\mathrm{g}
    = 0.511\ \mathrm{MeV}
    = (3.862 \times 10^{-11}\,\mathrm{cm})^{-1} .
\end{equation}
其他一些有用的数值和换算因子列在附录中。

\section*{相对论与张量}

我们关于相对论的约定遵循 Jackson（1975）、Bjorken 和 Drell（1964、1965）以及几乎所有近期的场论教科书。我们使用的度规张量为
\begin{equation}
  g_{\mu\nu} = \begin{pmatrix}
    1 & 0 & 0 & 0\\
    0 & -1 & 0 & 0\\
    0 & 0 & -1 & 0\\
    0 & 0 & 0 & -1
  \end{pmatrix},
\end{equation}
其中希腊指标取遍 $0,1,2,3$ 或 $t,x,y,z$。罗马指标——$i,j$ 等——只表示三个空间分量。重复指标在一切情形下都要求和。四矢量与普通数一样用浅斜体表示；三矢量用粗体表示；单位三矢量用带帽子的浅斜体字母表示。例如，
\begin{align}
  x^\mu &= (x^0, \mathbf{x}),\qquad x_\mu = g_{\mu\nu} x^\nu = (x^0, -\mathbf{x});\\
  p\cdot x &= g_{\mu\nu} p^\mu x^\nu = p^0 x^0 - \mathbf{p}\cdot\mathbf{x}.
\end{align}
一个有质量的粒子满足
\begin{equation}
  p^2 = p_\mu p^\mu = E^2 - |\mathbf{p}|^2 = m^2 .
\end{equation}
注意位移矢量 $x^\mu$ ``天然地取上升指标''，而导数算符
\begin{equation}
  \partial_\mu = \Bigl( \frac{\partial}{\partial t},\ \nabla \Bigr),
\end{equation}
则``天然地取下降指标''。

我们定义全反对称张量 $\epsilon^{\mu\nu\rho\sigma}$，使得
\begin{equation}
  \epsilon^{0123} = +1 .
\end{equation}
要当心，因为这蕴含 $\epsilon_{0123} = -1$ 和 $\epsilon^{1230} = -1$。（这一约定与 Jackson 一致，但与 Bjorken 和 Drell 不同。）

\section*{量子力学}

我们经常要用到单个量子力学粒子的薛定谔波函数。按照通常的约定，作用在这样的波函数上的能量和动量算符表示为：
\begin{equation}
  E = i\frac{\partial}{\partial t},\qquad \mathbf{p} = -i\nabla .
\end{equation}
这些方程可以合并为
\begin{equation}
  p^\mu = i\partial^\mu ,
\end{equation}
对 $\partial_\mu$ 上升指标恰好说明了负号的来源。平面波 $e^{-ik\cdot x}$ 具有动量
\begin{equation}
  k^\mu ,
\end{equation}
因为
\begin{equation}
  i\partial^\mu \bigl( e^{-ik\cdot x} \bigr) = k^\mu e^{-ik\cdot x} .
\end{equation}
记号``h.c.''表示厄米共轭。

量子力学中对自旋的讨论要用到泡利 $\sigma$ 矩阵：
\begin{equation}
  \sigma^1 = \begin{pmatrix}0&1\\1&0\end{pmatrix},\qquad
  \sigma^2 = \begin{pmatrix}0&-i\\i&0\end{pmatrix},\qquad
  \sigma^3 = \begin{pmatrix}1&0\\0&-1\end{pmatrix}.
\end{equation}
这些矩阵的乘积满足恒等式
\begin{equation}
  \sigma^i \sigma^j = \delta^{ij} + i\epsilon^{ijk} \sigma^k .
\end{equation}
定义线性组合 $\sigma^{\pm} = \frac{1}{2}(\sigma^1 \pm i\sigma^2)$ 是方便的；于是
\begin{equation}
  [\sigma^3, \sigma^\pm] = \pm 2 \sigma^\pm .
\end{equation}

\section*{傅里叶变换与分布}

我们将经常用到海维赛德阶跃函数 $\theta(x)$ 和狄拉克德尔塔函数 $\delta(x)$，它们的定义如下：
\begin{equation}
  \theta(x) = \begin{cases}1 & x>0\\ 0 & x<0 \end{cases},\qquad
  \int_{-\infty}^{\infty} \!\! dx\ \delta(x) = 1 .
\end{equation}
$n$ 维德尔塔函数记为 $\delta^{(n)}(\mathbf{x})$，它在除 $\mathbf{x}=0$ 之外的任何地方都为零，且满足
\begin{equation}
  \int d^n x\ \delta^{(n)}(\mathbf{x}) = 1 .
\end{equation}
在傅里叶变换中，因子 $2\pi$ 总是与动量积分一起出现。例如，在四维情形，
\begin{equation}
  f(x) = \int \frac{d^4k}{(2\pi)^4}\; e^{-ik\cdot x} f(k),\qquad
  f(k) = \int d^4x\; e^{ik\cdot x} f(x) .
\end{equation}
（在三维变换中，指数中的符号将分别是 $+$ 和 $-$。）当不会引起混淆时，$f(k)$ 上的波浪号有时会被省略。另一个需要记住的重要 $2\pi$ 因子出现在恒等式
\begin{equation}
  \int d^4x\; e^{ik\cdot x} = (2\pi)^4 \delta^{(4)}(k) .
\end{equation}

\section*{电动力学}

我们采用海维赛德-洛伦兹约定，在这种约定下因子 $4\pi$ 出现在库仑定律和精细结构常数中，而不是出现在麦克斯韦方程中。于是点电荷 $Q$ 的库仑势为
\begin{equation}
  \Phi = \frac{Q}{4\pi r} ,
\end{equation}
精细结构常数为
\begin{equation}
  \alpha = \frac{e^2}{4\pi} = \frac{e^2}{4\pi \hbar c} = \frac{1}{137} .
\end{equation}
符号 $e$ 表示电子的电荷，是一个负量（尽管其符号很少产生影响）。我们通常使用麦克斯韦方程的相对论形式：
\begin{equation}
  \partial_\mu F^{\mu\nu} = e J^\nu ,
\end{equation}
\begin{equation}
  \epsilon^{\mu\nu\rho\sigma}\ \partial_\nu F_{\rho\sigma} = 0 ,
\end{equation}
其中
\begin{equation}
  A^\mu = (\Phi,\ \mathbf{A}),\qquad
  F^{\mu\nu} = \partial^\mu A^\nu - \partial^\nu A^\mu
\end{equation}
而我们把 $e$ 从四矢量流密度中提取了出来
\begin{equation}
  J^\mu = (\rho,\ \mathbf{J}) .
\end{equation}

\section*{狄拉克方程}

我们的一些约定与 Bjorken 和 Drell（1964、1965）以及其他教科书不同：我们为狄拉克矩阵使用\emph{手征基}，为狄拉克旋量使用\emph{相对论归一化}。这些约定在第~3.2 节和第~3.3 节引入，并在附录中加以总结。

% ======================================================================
\chapter*{编者前言}
\addcontentsline{toc}{chapter}{编者前言}

如何以连贯的方式交流物理学中最激动人心、最活跃领域的最新进展，这个问题至今仍伴随着我们。物理学家数量的急剧增长，已使那些惯常的交流渠道变得大不如前。某一领域的专家要想跟上当前的文献已越来越困难；新手则只能感到困惑。人们所需要的，既是对某一领域前后一致的阐述，也是关于它的一种明确``观点''的呈现。在一个快速发展的领域，正式专著无法满足这样的需求，而综述文章似乎又已失宠。事实上，那些最积极地投身于发展某一领域的人，往往是最不愿意就此长篇大论的人。

\textit{Frontiers in Physics}（《物理学前沿》）丛书于 1961 年构想出来，力图从几个方面改善这种状况。一流的物理学家经常在他们特别感兴趣的领域开设一系列讲座、研究生研讨班或研究生课程。这样的讲座用于总结一个快速发展的领域的现状，而且很可能是当时唯一可用的连贯阐述。通常，讲座笔记（由讲课者、研究生或博士后准备）是存在的，并以复印形式在小范围内分发。\textit{Frontiers in Physics} 丛书的主要目的之一，就是让更广大的物理学家读者能够看到这些笔记。

随着 \textit{Frontiers in Physics} 的发展，第二类书——非正式教材/专著，介于讲座笔记与正式教材或专著之间的中间形态——在该丛书中扮演着越来越重要的角色。在非正式教材或专著中，作者对其讲座笔记作了加工，使书稿达到能够连贯总结一个新发展的领域——并附有参考文献和习题——的程度，既适合课堂教学，也适合个人自学。

在过去二十年里，无论在量子场论的概念框架方面，还是在其对凝聚态物理和基本粒子物理的应用方面，都取得了重大进展。鉴于量子场论的学习已成为物理专业研究生教育不可或缺的一部分，一本既能把这些新进展带给初学者、又不忽视基本概念的教科书，是极为可贵的。Michael Peskin 和 Daniel Schroeder 写的正是这样一本书，它以清晰的笔触描述了量子电动力学、重整化和非阿贝尔规范理论，同时让读者领略到即将展开的内容。因此，把这本非常精美的教材/专著收入 \textit{Frontiers in Physics} 丛书是再合适不过的了，我很高兴欢迎他们加入丛书作者的行列。

\vspace{2em}
\noindent 科罗拉多州阿斯彭\\
1995 年 8 月

\vspace{1em}
\noindent David Pines
